\documentclass[11pt]{article}

\usepackage[margin=1in]{geometry}
\usepackage{amsmath,amssymb}
\usepackage{hyperref}
\usepackage{booktabs}

\title{Bounding the Stable Time Step for Explicit Time Integration:\\
Stiffness, Linear Stability Theory, and Application to a\\
Penalty Spring--Dashpot Contact Model}
\author{}
\date{}

\begin{document}
\maketitle

\section{Motivation}

The particle simulation in this repository integrates a system of ODEs
\begin{equation}
  \dot{\mathbf{y}} = f(\mathbf{y}), \qquad \mathbf{y} = (x, y, v_x, v_y) \text{ per particle},
\end{equation}
using either semi-implicit (symplectic) Euler or a fixed-point (Picard) backward Euler
scheme, with forces dominated by a linear penalty spring plus a linear dashpot at each
active contact. \texttt{stability\_and\_accuracy.cpp} currently determines a stable
\texttt{dt} \emph{empirically}, by sweeping \texttt{dt} $\times$ \texttt{max\_force} and
classifying each combination as \texttt{UNSTABLE}, \texttt{STEADY}, or \texttt{TIMED\_OUT}.
This document derives the corresponding \emph{analytical} bound, so that the empirical
sweep can be understood as measuring a quantity that is, for this particular force model,
solvable in closed form.

\section{Linear stability theory}

\subsection{The Dahlquist test equation and regions of absolute stability}

For a general nonlinear system $\dot{\mathbf{y}} = f(\mathbf{y})$, freeze the state at
$\mathbf{y}_0$ and linearize:
\begin{equation}
  \dot{\mathbf{y}} \approx J \mathbf{y}, \qquad J = \left.\frac{\partial f}{\partial \mathbf{y}}\right|_{\mathbf{y}_0}.
\end{equation}
Each eigenvalue $\lambda$ of $J$ behaves, locally, like an independent copy of the scalar
\emph{Dahlquist test equation}
\begin{equation}
  \dot y = \lambda y, \qquad \lambda \in \mathbb{C},
  \label{eq:dahlquist}
\end{equation}
which is the basis of essentially all step-size stability theory for one-step methods
\cite{dahlquist1963}. A method's \emph{region of absolute stability} $S \subset \mathbb{C}$
is the set of values $z = h\lambda$ for which the method's amplification factor $R(z)$
(defined by $y_{n+1} = R(z)\, y_n$ applied to \eqref{eq:dahlquist}) satisfies $|R(z)| \le 1$.
A step size $h$ is linearly stable for a given $J$ iff $h\lambda \in S$ for every eigenvalue
$\lambda$ of $J$.

For \textbf{forward (explicit) Euler}, $y_{n+1} = y_n + h\lambda y_n = (1+h\lambda) y_n$, so
$R(z) = 1+z$ and
\begin{equation}
  S_{\text{FE}} = \{\, z \in \mathbb{C} : |1+z| \le 1 \,\},
\end{equation}
the closed disk of radius 1 centered at $-1$. For real $\lambda < 0$ this reduces to the
familiar
\begin{equation}
  h \le \frac{2}{|\lambda|}.
  \label{eq:fe-real-bound}
\end{equation}

For \textbf{backward Euler}, $y_{n+1} = y_n + h\lambda y_{n+1}$, so $R(z) = 1/(1-z)$ and
$S_{\text{BE}} \supseteq$ the entire left half-plane $\mathrm{Re}(z) \le 0$ — backward Euler
is \emph{A-stable} \cite{dahlquist1963}: linearly unconditionally stable for any $\lambda$
with $\mathrm{Re}(\lambda) \le 0$, regardless of $h$.

\subsection{Stiffness}

Given the eigenvalues $\{\lambda_i\}$ of $J$ relevant to the dynamics of interest, the
\emph{stiffness ratio} is
\begin{equation}
  S = \frac{\max_i |\mathrm{Re}(\lambda_i)|}{\min_i |\mathrm{Re}(\lambda_i)|}.
\end{equation}
A system is called \emph{stiff} when $S \gg 1$: an explicit method is forced by
\eqref{eq:fe-real-bound} to resolve the fastest (often physically uninteresting) mode just
to remain stable, even though the mode of actual interest decays or oscillates far more
slowly. This is the definition introduced by Curtiss and Hirschfelder in the paper that
coined the term "stiff equations," originally for chemical-kinetics ODEs
\cite{curtisshirschfelder1952}.

\subsection{Backward Euler with Picard iteration: a separate constraint}

A-stability of backward Euler is a statement about the \emph{linear} amplification factor;
it says nothing about whether the nonlinear equation $\mathbf{y}_{n+1} = \mathbf{y}_n + h
f(\mathbf{y}_{n+1})$ can actually be solved by fixed-point (Picard) iteration, as
\texttt{BackwardEulerPicardKernel} does. Writing the iteration as
$\mathbf{y}^{(k+1)} = \mathbf{y}_n + h f(\mathbf{y}^{(k)})$, this is a contraction (and
therefore converges, by the Banach fixed-point theorem) only if
\begin{equation}
  h \left\| \frac{\partial f}{\partial \mathbf{y}} \right\| < 1,
  \label{eq:picard-contraction}
\end{equation}
i.e.\ $h$ times the local Lipschitz constant of $f$ must be less than 1. This is a
genuinely different — usually much looser — bound than forward Euler's linear stability
bound \eqref{eq:fe-real-bound}, and it constrains \emph{iteration convergence}, not the
stability of the underlying scheme.

\subsection{Estimating the spectral radius numerically}

For large or expensive-to-differentiate systems, one rarely forms and diagonalizes $J$
explicitly. Two standard cheap alternatives:

\begin{itemize}
  \item \textbf{Gershgorin circle theorem}: every eigenvalue of $J$ lies in some disk
    centered at $J_{ii}$ with radius $\sum_{j \ne i} |J_{ij}|$ — a bound computable in a
    single pass over $J$ with no eigensolve.
  \item \textbf{Power iteration}, applied either to $J$ directly or to a finite-difference
    directional derivative $(f(\mathbf{y}+\varepsilon \mathbf{v}) - f(\mathbf{y}))/\varepsilon$
    so that $J$ is never assembled at all. This is exactly the technique used by explicit
    stabilized Runge--Kutta--Chebyshev (RKC) integrators to adaptively re-estimate the
    spectral radius of $J$ every step \cite{sommeijershampineverwer1997}.
\end{itemize}

The \texttt{dt} $\times$ \texttt{max\_force} sweep in \texttt{stability\_and\_accuracy.cpp}
is, in this sense, a brute-force empirical stand-in for the same question these methods
answer analytically: it \emph{measures} the stability boundary by bisection/sweep rather
than predicting it from $J$.

\section{Closed-form bound for the penalty spring--dashpot contact model}

\subsection{The linearized contact ODE}

Each active contact in \texttt{compute\_rhs\_kernel.cu} applies a linear penalty spring
plus a linear dashpot along the contact normal. Restricting attention to the
one-dimensional normal direction of a single contact, with overlap $x$ (positive when
penetrating), this is exactly a damped harmonic oscillator:
\begin{equation}
  m_{\text{eff}}\, \ddot x + c\, \dot x + k\, x = 0,
  \label{eq:sho}
\end{equation}
where, matching the code's own quantities,
\begin{align}
  k &= \frac{\texttt{max\_force}}{\texttt{particle\_radius}}, \\
  m_{\text{eff}} &= \begin{cases}
    m_i & \text{particle--wall/floor/ceiling (infinite-mass boundary limit)} \\
    \dfrac{m_i m_j}{m_i + m_j} & \text{particle--particle (reduced mass)}
  \end{cases}
\end{align}
and $c$ is given by \texttt{restitution\_to\_damping()}, i.e.\ derived from the target
coefficient of restitution $e \in [0,1]$ via the standard DEM restitution-to-damping
relation \cite{tsuji1992,direnzodimaio2004}.

Define the undamped natural frequency and damping ratio
\begin{equation}
  \omega_n = \sqrt{\frac{k}{m_{\text{eff}}}}, \qquad
  \zeta = \frac{c}{2\sqrt{k\, m_{\text{eff}}}}.
\end{equation}

\subsection{Eigenvalues of the linearized contact}

Writing \eqref{eq:sho} as a first-order system with $y_1 = x$, $y_2 = \dot x$:
\begin{equation}
  \frac{d}{dt}\begin{pmatrix} y_1 \\ y_2 \end{pmatrix}
  = \underbrace{\begin{pmatrix} 0 & 1 \\ -\omega_n^2 & -2\zeta\omega_n \end{pmatrix}}_{J}
  \begin{pmatrix} y_1 \\ y_2 \end{pmatrix},
\end{equation}
with characteristic polynomial $\lambda^2 + 2\zeta\omega_n \lambda + \omega_n^2 = 0$, giving
\begin{equation}
  \lambda = -\zeta\omega_n \pm \omega_n\sqrt{\zeta^2 - 1}.
  \label{eq:eigs}
\end{equation}
For $0 \le \zeta < 1$ (underdamped — i.e.\ $e > 0$, the usual case) the eigenvalues are the
complex-conjugate pair
\begin{equation}
  \lambda_{\pm} = -\zeta\omega_n \pm i\,\omega_n\sqrt{1-\zeta^2}.
  \label{eq:eigs-complex}
\end{equation}

\subsection{Forward Euler stability bound on the damped oscillator}

Apply the forward Euler stability condition $|1 + h\lambda| \le 1$ to \eqref{eq:eigs-complex}.
Write $z = 1 + h\lambda_+ = (1 - h\zeta\omega_n) + i\, h\omega_n\sqrt{1-\zeta^2}$. Then
\begin{align}
  |z|^2 &= (1 - h\zeta\omega_n)^2 + h^2\omega_n^2(1-\zeta^2) \\
        &= 1 - 2h\zeta\omega_n + h^2\zeta^2\omega_n^2 + h^2\omega_n^2 - h^2\omega_n^2\zeta^2 \\
        &= 1 - 2h\zeta\omega_n + h^2\omega_n^2.
\end{align}
Requiring $|z|^2 \le 1$:
\begin{equation}
  -2h\zeta\omega_n + h^2\omega_n^2 \le 0
  \quad\Longleftrightarrow\quad
  h\left(h\omega_n^2 - 2\zeta\omega_n\right) \le 0
  \quad\Longleftrightarrow\quad
  h \le \frac{2\zeta}{\omega_n} \qquad (h > 0).
  \label{eq:fe-bound-damped}
\end{equation}

\textbf{Consequence.} As $\zeta \to 0$ (undamped, perfectly elastic contact, $e \to 1$),
the bound \eqref{eq:fe-bound-damped} goes to \emph{zero}. This recovers the well-known fact
that plain forward Euler is \emph{unconditionally unstable} for an undamped oscillator, for
any step size $h > 0$ — the eigenvalues in \eqref{eq:eigs-complex} sit exactly on the
imaginary axis when $\zeta=0$, and the forward Euler stability disk $S_{\text{FE}}$ touches
the imaginary axis only at the origin. Damping is the \emph{only} thing buying a nonzero
stable step for plain forward Euler on this model.

\subsection{Semi-implicit (symplectic) Euler}

The scheme actually used by \texttt{SemiImplicitEulerKernel} updates velocity first, then
position using the \emph{updated} velocity:
\begin{align}
  v_{n+1} &= v_n + h\big(-\omega_n^2 x_n - 2\zeta\omega_n v_n\big), \\
  x_{n+1} &= x_n + h\, v_{n+1}.
\end{align}
In matrix form, $\begin{pmatrix} x_{n+1} \\ v_{n+1} \end{pmatrix} = M \begin{pmatrix} x_n \\ v_n \end{pmatrix}$ with
\begin{equation}
  M = \begin{pmatrix} 1 - h^2\omega_n^2 & h - 2\zeta\omega_n h^2 \\ -h\omega_n^2 & 1 - 2\zeta\omega_n h \end{pmatrix}.
\end{equation}
In the undamped limit $\zeta = 0$,
\begin{equation}
  M = \begin{pmatrix} 1 - h^2\omega_n^2 & h \\ -h\omega_n^2 & 1 \end{pmatrix},
  \qquad \det M = 1, \qquad \operatorname{tr} M = 2 - h^2\omega_n^2.
\end{equation}
Since $\det M = 1$, the eigenvalues of $M$ are reciprocal ($\mu_1\mu_2 = 1$); they lie on
the unit circle (marginal, bounded stability) exactly when $|\operatorname{tr} M| \le 2$,
i.e.
\begin{equation}
  \Big| 2 - h^2\omega_n^2 \Big| \le 2 \quad\Longleftrightarrow\quad h \le \frac{2}{\omega_n}.
  \label{eq:si-bound}
\end{equation}
Unlike \eqref{eq:fe-bound-damped}, this bound is \emph{finite even as $\zeta \to 0$}: this
is the standard reason symplectic/leapfrog-type schemes are used for oscillatory or
near-elastic contact dynamics rather than plain forward Euler, and it is exactly the
$h \lesssim 2/\omega_n = 2\sqrt{m_{\text{eff}}/k}$ scaling that appears throughout the DEM
literature as the \emph{critical} or \emph{Rayleigh} time step
\cite{cundallstrack1979,osullivanbray2004}.

\subsection{Practical bound for this simulation}

Since $k$, \texttt{particle\_radius}, and per-material mass are static configuration, the
binding constraint is set by whichever \emph{currently active} contact has the largest
$\omega_n = \sqrt{k/m_{\text{eff}}}$ — i.e.\ the lightest, stiffest, most nearly-elastic
pair, which in this codebase's material set is almost always a wall contact using a light
particle's own mass (the infinite-mass boundary limit already used in
\texttt{restitution\_to\_damping()}), or a snow--snow contact if snow is the lightest
material. Because this does not depend on the runtime particle configuration, it can be
computed once from \texttt{config.yaml} — via \eqref{eq:si-bound} for the semi-implicit
integrator, or empirically applying a conservative safety factor (0.2--0.4, per
\cite{osullivanbray2004}) — as a much cheaper substitute for, or a sanity check on, the
brute-force \texttt{dt} sweep in \texttt{stability\_and\_accuracy.cpp}.

\section{Summary}

\begin{itemize}
  \item Local linear stability of any explicit one-step method reduces to checking each
    eigenvalue of the local Jacobian against the method's region of absolute stability
    (Section 2.1).
  \item Stiffness is the ratio of fastest to slowest relevant eigenvalue magnitudes
    (Section 2.2); it is what forces an explicit step size far below what the dynamics of
    interest would otherwise require.
  \item For large systems, the spectral radius of $J$ can be estimated cheaply via
    Gershgorin's theorem or power iteration, without ever forming $J$ explicitly
    (Section 2.4) — the same idea that underlies adaptive RKC-family integrators.
  \item This simulation's contact model is, per-contact, a damped harmonic oscillator, so
    its stable step size is solvable in closed form: plain forward Euler is unconditionally
    unstable in the undamped limit (\eqref{eq:fe-bound-damped}), while the semi-implicit
    Euler scheme actually used has the finite bound $h \le 2/\omega_n$
    (\eqref{eq:si-bound}), matching the classic Rayleigh/critical time step result from the
    DEM literature.
\end{itemize}

\begin{thebibliography}{9}

\bibitem{dahlquist1963}
G.~Dahlquist,
\emph{A special stability problem for linear multistep methods},
BIT Numerical Mathematics \textbf{3}, 27--43 (1963).

\bibitem{curtisshirschfelder1952}
C.~F.~Curtiss and J.~O.~Hirschfelder,
\emph{Integration of stiff equations},
Proceedings of the National Academy of Sciences \textbf{38}(3), 235--243 (1952).

\bibitem{sommeijershampineverwer1997}
B.~P.~Sommeijer, L.~F.~Shampine, and J.~G.~Verwer,
\emph{RKC: an explicit solver for parabolic PDEs},
Journal of Computational and Applied Mathematics \textbf{88}(2), 315--326 (1997).

\bibitem{cundallstrack1979}
P.~A.~Cundall and O.~D.~L.~Strack,
\emph{A discrete numerical model for granular assemblies},
G\'{e}otechnique \textbf{29}(1), 47--65 (1979).

\bibitem{osullivanbray2004}
C.~O'Sullivan and J.~D.~Bray,
\emph{Selecting a suitable time step for discrete element simulations that use the central
difference time integration scheme},
Engineering Computations \textbf{21}(2/3/4), 278--303 (2004).

\bibitem{tsuji1992}
Y.~Tsuji, T.~Tanaka, and T.~Ishida,
\emph{Lagrangian numerical simulation of plug flow of cohesionless particles in a
horizontal pipe},
Powder Technology \textbf{71}(3), 239--250 (1992).

\bibitem{direnzodimaio2004}
A.~Di Renzo and F.~P.~Di Maio,
\emph{Comparison of contact-force models for the simulation of collisions in
DEM-based granular flow codes},
Chemical Engineering Science \textbf{59}(3), 525--541 (2004).

\bibitem{gear1971}
C.~W.~Gear,
\emph{Numerical Initial Value Problems in Ordinary Differential Equations},
Prentice-Hall (1971).

\bibitem{hairerwanner1996}
E.~Hairer and G.~Wanner,
\emph{Solving Ordinary Differential Equations II: Stiff and Differential-Algebraic
Problems}, 2nd ed., Springer Series in Computational Mathematics (1996).

\bibitem{courantfriedrichslewy1928}
R.~Courant, K.~Friedrichs, and H.~Lewy,
\emph{\"Uber die partiellen Differenzengleichungen der mathematischen Physik},
Mathematische Annalen \textbf{100}(1), 32--74 (1928).

\end{thebibliography}

\end{document}
